\chapter{Optimización en dimensión finita}
\noindent
La única experiencia que hemos tenido en la carrera con el mundo de la optimización es la optimización de funciones reales de variable real. Como el lector pueda llegar a comprender este punto de partida es demasiado básico para el nivel de la asignatura, que comienza explicnado la optimización en dimensión infinita.\\

\noindent
Por ello, es recomendable repasar el documento que le haya proporcionado su profesor sobre conceptos clave que hay que tener en cuenta a la hora de realizar optimización en varias dimensiones, de forma pormenorizada. En este capítulo recuperamos (o al menos mencionamos) los conceptos más importantes que hemos de tener en cuenta después de esta pasada rápida por el mundo de la optimización en dimensión finita, para que el choque con la optimización en dimensión finita no nos sea tan grande como pueda parecer.\\

\noindent
Para el desarrollo de este tema, conviene tener en cuenta a la hora de resolver un problema de optimización en dimensión finita: % // TODO: Quizás desarrollar más este punto, o no
\begin{itemize}
    \item Análisis de derivadas primeras y segundas: formas cuadráticas y polinomio de Taylor.
    \item Complicaciones en la frontera.
    \item Optimización con restricciones de igualdad\footnote{En el caso de las restricciones de tipo desigualdad tenemos los multiplicadores KKT (Karush-Kuhn-Tucker).}, método de los multiplicadores de Lagrange.
    \item Nociones de convexidad permiten pasar de nociones locales a globales.
    \item Derivadas direccionales.
    \item Programación lineal\footnote{Hay problemas similares, conocidos como programación cuadrática y programación convexa.}.
    \item Teorema de Weierstrass para la existencia de óptimos, previa a buscarlos.
\end{itemize}

\begin{teo}[de Weierstrass]
    Sea $K\subset \mathbb{R}^d$ un conjunto compacto y $f:K\to \mathbb{R}$ una aplicación continua. Entonces $f$ alcanza su mínimo.
    \begin{proof}
        Tomando $\alpha = \inf\limits_{x\in K}f(x)$, sabemos que existe una sucesión $\{x_n\}$ de puntos de $K$ de forma que $\{f(x_n)\}\to \alpha$. Por compacidad ha de existir $x\in K$ de forma que $\{x_n\}\to x$, y por continuidad de $f$ tenemos que $\{f(x_n)\}\to f(x)$, de donde $f(x) = \alpha$.
    \end{proof}
\end{teo}

\begin{definicion}[Semicontinua inferiormente]
    Sea $X$ un espacio topológico 1AN, decimos que una aplicación $f:X\to \mathbb{R}$ es semicontinua inferiormente si para toda sucesión convergente $\{x_n\}$ de puntos de $X$ se tiene que:
    \begin{equation*}
        \liminf \{f(x_n)\} \geq f\left(\lim \{x_n\}\right)
    \end{equation*}
\end{definicion}

\begin{teo}[Método directo del cálculo de variaciones]
    Sea $K$ un espacio topológico 1AN compacto y $F:K\to \mathbb{R}$ una aplicación semicontinua inferiormente. Entonces $F$ alcanza su mínimo.
    \begin{proof}
        Tomando $\alpha = \inf\limits_{x\in K}F(x)$, sabemos que existe una sucesión $\{x_n\}$ de puntos de $K$ de forma que $\{F(x_n)\}\to \alpha$. Por compacidad ha de existir $x\in K$ de forma que $\{x_n\}\to x$. Tenemos así que:
        \begin{equation*}
            \alpha = \lim \{F(x_n)\} = \liminf \{F(x_n)\} \geq F(x)
        \end{equation*}
        Y como $\alpha = \inf\limits_{x\in K}F(x)$ tendrá que ser $F(x) = \alpha$.
    \end{proof}
\end{teo}

\chapter{Optimización en dimensión infinita}

\section{Derivada de Gateaux}
\noindent
En dimensión infinita tenemos también el concepto de derivadas direccionales.\\

\noindent
Sea $X$ un espacio normado, $D\subset X$ y sea $F:D\to \mathbb{R}$. Fijamos $y\in D$, diremos que $\varphi \in X\setminus \{0\}$ es \textbf{dirección admisible para $y$ en $D$} si existe $\varepsilon_0>0$ tal que $y+\varepsilon\varphi \in D$ para\footnote{Podríamos definirlo en $\left]0,\varepsilon_0\right[$, pero por motivmos pedagógicos nos restringiremos a este caso, por simplicidad a la hora de dar los enunciados.} $\varepsilon \in \left]-\varepsilon_0,\varepsilon_0\right[$. Por tanto, está definida la correspondencia:
\begin{equation*}
    \varepsilon\longmapsto F(y+\varepsilon\varphi), \qquad \varepsilon \in \left]-\varepsilon_0,\varepsilon_0\right[
\end{equation*}
Supongamos que es derivable, en cuyo caso diremos que $F$ es derivable en el sentido de Gateaux en $y$ a lo largo de $\varphi$. La derivada correspondiente se define como:
\begin{equation*}
    \delta_\varphi F(y) = \lim_{\varepsilon\to0} \frac{F(y+\varepsilon\varphi)-F(y)}{\varepsilon} = \left[\frac{d}{d\varepsilon}F(y+\varepsilon\varphi)\right]_{\big|\varepsilon=0}
\end{equation*}
Diremos que $F$ es derivable en el sentido de Gateaux en un punto $y\in D$ si $F$ es derivable en el sentido de Gateaux en $y$ para toda dirección admisible $\varphi$.

\noindent
Diremos que $y_\ast\in D$ es un \textbf{punto crítico} (o estacionario) de $F$ si\footnote{También incluimos aquí los puntos en los que no hay ninguna dirección admisible.}:
\begin{equation*}
    \delta_\varphi F(y_\ast) = 0 \qquad \forall \varphi \text{\ dirección admisible}
\end{equation*}

\begin{ejemplo}
    Consideremos $D = \{C^1[-1,1] : y(-1) = 1 = y(1)\}\subset C^1[-1,1]$. Para $y(x) = x^2 \in D$, ¿son direcciones admisibles las siguientes?
    \begin{align*}
        \varphi_1(x) &= x^2 \\
        \varphi_2(x) &= 1-x^2
    \end{align*}
    Para que sean direcciones admisibles necesitamos comprobar que
    \begin{equation*}
        y_{\varepsilon}(x) = y(x) + \varepsilon\varphi_i(x)\in D \qquad \text{para $\varepsilon \in \left]-\varepsilon_0,\varepsilon_0\right[$}, \quad i \in \{1,2\}
    \end{equation*}
    \begin{enumerate}
        \item Para $y_\varepsilon(x) = x^2+\varepsilon x^2 = x^2(1+\varepsilon)$, tenemos que $y_{\varepsilon}(-1) = 1+\varepsilon\neq 1$, por lo que no es una dirección admisible para ningún $\varepsilon$.
        \item Para $y_{\varepsilon}(x) = x^2+\varepsilon (1-x^2)$ tenemos que:
            \begin{align*}
                y_{\varepsilon}(-1) &= 1 + \varepsilon(1-1) = 1 \\
                y_\varepsilon(1) &= 1 + \varepsilon(1-1) = 1
            \end{align*}
            Por lo que sí que es una dirección admisible.
    \end{enumerate}
    Cualquier $\varphi \in C^1[-1,1]$ con $\varphi(-1) = 0 = \varphi(1)$ será una dirección admisible.
\end{ejemplo}

\begin{notacion}
    Notaremos usualmante:
    \begin{equation*}
        C^1_0([a,b]) = \{y\in C^1([a,b]) : y(a) = 0 = y(b)\}
    \end{equation*}
\end{notacion}

\begin{prop}
    Sea $X$ un espacio normado, sea $D\subset X$ y $F:D\to \mathbb{R}$. Si $y_0\in D$ es extremo relativo de $F$ y $F$ es derivable de Gateaux en $y_0$, entonces $y_0$ es un punto crítico.
    \begin{proof}
        Supongamos que en $y_0$ se alcanza un mínimo local (en caso de máximo es análogo), es decir, existe un entorno de $y_0$ en el cual $F(y_0)\leq F(y)$ para cualquier $y$ en dicho entorno. Sin pérdida de generalidad, podemos suponer que es de la forma $D\cap B(y_0,\delta)$, para $\delta\in \mathbb{R}^+$.\\

        \noindent
        Dada $\varphi \in X$ dirección admisible para $y_0$, definimos $y_{\varepsilon} = y_0 + \varepsilon\varphi$, con lo que existe $\varepsilon_0>0$ tal que $y_\varepsilon \in D\quad \forall \varepsilon \in \left]-\varepsilon_0,\varepsilon_0\right[$. Observemos que:
        \begin{equation*}
            \|y_0 - y_{\varepsilon}\| = |\varepsilon| \|\varphi\| \quad\Longrightarrow\quad y_{\varepsilon}\in B(y_0,\delta) \quad \forall \varepsilon \in \left]\frac{-\delta}{\|\varphi\|}, \frac{\delta}{\|\varphi\|}\right[
        \end{equation*}
        Sea $\overline{\varepsilon} = \min\left\{\varepsilon_0,\frac{\delta}{\varphi}\right\}$, la correspondencia $\varepsilon\longmapsto F(y_{\varepsilon})$ está definida en $\left]-\overline{\varepsilon},\overline{\varepsilon}\right[$, es derivable y tiene un mínimo local en $\varepsilon=0$. En dicho caso, como $F$ es derivable de Gateaux en $y_0$ respecto a $\varphi$, obtenemos que $\delta_\varphi F(y_0) = 0$. Como la dirección admisible era arbitraria, se cumple para cualquiera.
    \end{proof}
\end{prop}

\section{Estudio de funcionales integrales}
\noindent
Tomaremos $X = C^1(I)$, con $I=[x_0,x_1]$ para $x_0< x_1$. Usaremos la norma:
\begin{equation*}
    \|y\|_{C^1} = \|y\|_\infty + \|y'\|_\infty
\end{equation*}
Consideraremos funcionales dados por: 
\begin{equation*}
    \cc{F}(y) = \int_{x_0}^{x_1} F(x,y(x),y'(x))~dx 
\end{equation*}
Típicamente consideraremos\footnote{O quizás en otro dominio de definición.} $F:I\times\mathbb{R}\times\mathbb{R}\to \mathbb{R}$ integrable, e indicaremos sus variables por $F = F(x,y,p)$. 

\begin{notacion}
    Usaremos los convenios:
    \begin{equation*}
        F_x = \frac{\partial F}{\partial x}, \qquad F_y = \frac{\partial F}{\partial y}, \qquad F_p = \frac{\partial F}{\partial p}
    \end{equation*}
    Y los análogos para derivadas parciales segundas:
    \begin{equation*}
        F_{yp} = \frac{\partial^2 F}{\partial y\partial p}
    \end{equation*}
\end{notacion}

\begin{lema}[Fundamental del cálculo de variaciones]
    Sea $I = [x_0,x_1]$, sea $f\in C^1(I)$. Si:
    \begin{equation*}
        \int_{x_0}^{x_1} f(x)\varphi(x)~dx = 0  \qquad \forall \varphi \in C_0(I)
    \end{equation*}
    Entonces $f(x) = 0 \quad \forall x\in I$.
    \begin{proof}
        Por contrarrecíproco, podemos suponer que $\exists \overline{x}\in \left]x_0,x_1\right[$ tal que $f(\overline{x})>0$. Así, existe un entorno de $\overline{x}$ tal que $f(x) > 0 \quad \forall x\in \left]\overline{x_0},\overline{x_1}\right[$. Podemos encontrar $\varphi \in C_0(I)$ tal que
        \begin{equation*}
            \int_{x_0}^{x_1} f(x)\varphi(x)~dx  > 0
        \end{equation*}
        Por ejemplo, podemos construir $\varphi$ a partir de la función:
        \begin{equation*}
            \rho(x) = \left\{\begin{array}{ll}
                e^{\frac{1}{|x|^2-1}} & \text{si\ } |x| < 1 \\
                 0 & \text{en otro caso\ } 
            \end{array}\right. 
        \end{equation*}
        donde $\rho \in C_0^\infty([-1,1])$.
    \end{proof}
\end{lema}

\begin{observacion}
    Varios comentarios:
    \begin{itemize}
        \item En el Lema fundamental se puede cambiar $C_0(I)$ por $C_0^1(I), \ldots, C_0^k(I), \ldots, C_0^\infty(I)$.
        \item Hay una versión alternativa del Lema, que a continuación se enuncia.
        \item De hecho, en dicho enunciado podemos también exigir para toda $\varphi \in C_c^\infty(I)$.
    \end{itemize}
\end{observacion}

\begin{lema}
    Sea $I = \left]x_0,x_1\right[$, $f\in L^1_{\text{loc}}(I)$. Si:
    \begin{equation*}
        \int_{x_0}^{x_1} f(x)\varphi(x)~dx = 0 \qquad \forall \varphi \in C_c(I)
    \end{equation*}
    Entonces, $f(x) = 0$ para casi todo $x\in I$.
    \begin{proof}
        El dual de $C_c(I)$ contiene a $L^1_{\text{loc}}(I)$ y se extrae a partir de ahí.
    \end{proof}
\end{lema}

\begin{teo}\label{teo:principal}
    Sea $I=[x_0,x_1]$, $F\in C^2(I\times\mathbb{R}\times\mathbb{R})$. Consideramos el funcional
    \begin{equation*}
        \cc{F}(y) = \int_{x_0}^{x_1} F(x,y(x),y'(x))~dx 
    \end{equation*}

    definido sobre:
    \begin{equation*}
        D = \{y\in C^1([x_0,x_1]) : y(x_0) = y_0, y(x_1) = y_1\}, \quad y_0,y_1\in \mathbb{R}
    \end{equation*}
    Dado $y_\ast\in D\cap C^2(I)$ punto crítico de $\cc{F}$, satisface la ecuación de Euler-Lagrange:
    \begin{gather*}
        F_y(x,y_\ast(x),y_\ast'(x)) - \frac{d}{dx}\left[F_p(x,y_\ast(x),y_\ast'(x))\right] = 0 \qquad \forall x\in I \\
        \\
        \boxed{F_y - \frac{d}{dx}(F_p) = 0}
    \end{gather*}
    \begin{proof}
        Dado $y\in D$, $\varphi \in X\setminus \{0\}$, notaremos $y_\varepsilon = y+\varepsilon\varphi$, calculemos $\delta_\varphi\cc{F}(y)$ formalmente:
        \begin{gather*}
            \frac{d y_\varepsilon}{d \varepsilon} = \varphi \\
            y_\varepsilon' = y' + \varepsilon\varphi' \Longrightarrow \frac{d y'_\varepsilon}{d \varepsilon} = \varphi'
        \end{gather*}
        De donde:
        \begin{align*}
            \frac{d }{d \varepsilon} \cc{F}(y_\varepsilon) &= \frac{d }{d \varepsilon}\int_{x_0}^{x_1} F(x,y_\varepsilon(x),y_\varepsilon'(x))~dx  \\ 
                                                           &= \int_{x_0}^{x_1} F_y(x,y_\varepsilon(x),y_\varepsilon'(x))\frac{d y_\varepsilon}{d \varepsilon} + F_p(x,y_\varepsilon(x),y_\varepsilon'(x))\frac{dy_\varepsilon'}{d\varepsilon}~dx  \\
                                                           &= \int_{x_0}^{x_1} F_y(x,y_\varepsilon(x),y_\varepsilon'(x))\varphi(x) + F_p(x,y_\varepsilon(x),y_\varepsilon'(x))\varphi'(x)~dx 
        \end{align*}
        Por lo que:
        \begin{equation*}
            \left[\frac{d}{d\varepsilon}\cc{F}(y_\varepsilon)\right]_{\big|\varepsilon=0} = \int_{x_0}^{x_1} F_y(x,y(x),y'(x))\varphi(x) + F_p(x,y(x),y'(x))\varphi'(x)~dx 
        \end{equation*}
        Así:
        \begin{equation*}
            \delta_\varphi \cc{F}(y) = \int_{x_0}^{x_1} F_y\varphi + F_p\varphi'~dx 
        \end{equation*}
        Nos queremos quitar $\varphi'$ de esta expresión, lo conseguiremos integrando por partes:
        \begin{multline*}
            \int_{x_0}^{x_1} F_p(x,y(x),y'(x))\varphi'(x)~dx  = -\int_{x_0}^{x_1} \frac{d}{dx}[F_p(x,y(x),y'(x))]\varphi(x)~dx \\+ F_p(x_1,y(x_1),y'(x_1))\varphi(x_1) - F_p(x_0,y(x_0),y'(x_0))\varphi(x_0)
        \end{multline*}
        De donde:
        \begin{equation*}
            \delta_\varphi \cc{F}(y) = \int_{x_0}^{x_1} \left(F_y - \frac{d}{dx}F_p\right)\varphi~dx  + [F_p\varphi]_{x_0}^{x_1}
        \end{equation*}
        ¿Cuáles son las direcciones admisibles? Dada $y\in D$, ¿cómo ha de ser $\varphi$ para que $y_\varepsilon = y+\varepsilon\varphi \in D$? Tomando $\varphi \in C^1([x_0,x_1])$ tenemos la condición de regularidad, y vemos que:
        \begin{align*}
            y_\varepsilon(x_0) &= y(x_0) + \varepsilon\varphi(x_0) = y_0 \quad\Longleftrightarrow\quad \varphi(x_0) = 0
        \end{align*}
        Análogamente necesitamos que $\varphi(x_1) = 0$. Por tanto, la direcciones admisibles son $\varphi \in C^1_0([x_0,x_1])$. Usando estas direcciones admisibles, vemos que:
        \begin{equation*}
            \delta_\varphi \cc{F}(y) = \int_{x_0}^{x_1} \left(F_y - \frac{d}{dx}F_p\right)\varphi~dx  + \cancelto{0}{[F_p\varphi]_{x_0}^{x_1}}
        \end{equation*}
        Para cualquier dirección admisible $\varphi$. Usando que $y_\ast$ es punto crítico, tenemos:
        \begin{equation*}
            \delta_\varphi \cc{F}(y_\ast) = \int_{x_0}^{x_1}\left( F_y(x,y_\ast(x),y_\ast'(x)) - \frac{d}{dx}[F_p(x,y_\ast(x),y_\ast'(x))]\right)\varphi(x)~dx  = 0
        \end{equation*}
        Para toda $\varphi \in C_0^1(I)$. Por tanto, por el Lema fundamental del cálculo de variaciones tenemos la ecuación de Euler-Lagrange.
    \end{proof}
\end{teo}

\noindent
Si los candidatos no son de clase 2, podremos definir soluciones de la ecuación de Euler-Lagrange a partir de la fórmula:
\begin{equation*}
    \delta_\varphi\cc{F}(y) = \int_{x_0}^{x_1} F_y\varphi + F_p \varphi'~dx 
\end{equation*}
Pero involucra conceptos de análisis funcional más finos.

\begin{ejemplo}
    Estudiar las curvas planas que unen dos puntos dados con mínima longitud.\\

    \noindent
    Tomamos $A = (x_0,y_0), B = (x_1,y_1) \in \mathbb{R}^2$ y supondremos que $x_0<x_1$. Vamos a considerar únicamente aquellas curvas que pueden ser descritas como gráficas de una función\footnote{Pues nos parece que las curvas que optimizan el problema son de este tipo y simplifica mucho el problema.}, es decir, aplicaciones $x\longmapsto (x,y(x))$ con:
    \begin{equation*}
        y \in  D = \{f\in C^1([x_0,x_1]) : f(x_0) = y_0, f(x_1) = y_1\}
    \end{equation*}
    Así, las curvas tienen extremos en los puntos $A$ y $B$. Para este tipo de curvas se calcula su longitud como:
    \begin{equation*}
        L(y) = \int_{x_0}^{x_1} \|(1,y'(x))\|_2~dx = \int_{x_0}^{x_1} \sqrt{1+{(y'(x))}^{2}}~dx 
    \end{equation*}
    Así, tenemos $F= F(p) = \sqrt{1+p^2}$. La ecuación de Euler-Lagrange en este caso es:
    \begin{gather*}
        F_p = \frac{p}{\sqrt{1+p^2}} \\
        0 = F_y - \frac{d}{dx}(F_p) = -\frac{d}{dx} \left(\frac{y'(x)}{\sqrt{1+{(y'(x))}^{2}}}\right) \qquad \forall x\in [x_0,x_1]
    \end{gather*}
    Por lo que todos los puntos críticos $y(x)$ de $L$ satisfacen esta ecuación diferencial, de donde podemos deducir que para cierto $c_1\in \mathbb{R}$:
    \begin{equation*}
        \frac{y'(x)}{\sqrt{1+{(y'(x))}^{2}}} = c_1  \quad\Longrightarrow\quad y'(x) = c_1\sqrt{1+{(y'(x))}^{2}}
    \end{equation*}
    
    de donde:
    \begin{equation*}
        {(y'(x))}^{2} = c_1 \left( 1+{(y'(x))}^{2} \right) \quad\Longleftrightarrow\quad {(y'(x))}^{2} = c_2
    \end{equation*}
    para cierto $c_2\in \mathbb{R}$. Como $y\in C^1([x_0,x_1])$, tenemos entonces que $y'(x)$ es constante $\forall x\in [x_0,x_1]$, con lo que $y(x) = c_3 + c_4x\quad \forall x\in [x_0,x_1]$.\\

    \noindent
    Las constantes $c_3$ y $c_4$ se determinan por las condiciones $y(x_i) = y_i$, $i \in \{0,1\}$, obteniendo un único \underline{candidato a solución}, que se corresponde con la línea recta. Habría que comprobar que este candidato es efectivamente un mínimo de la función a optimizar.
\end{ejemplo}

\begin{definicion}
    Diremos que $y\in X$ es un \textbf{extremal} del funcional $\cc{F}$ si resuelve la ecuación de Euler-Lagrange asociada (notemos que no se pide $y\in D$).
\end{definicion}

\section{Problemas con frontera libre}
\noindent
En las condiciones del Teorema~\ref{teo:principal}, supongamos que 
\begin{equation*}
    D = \{y\in C^1([x_0,x_1]) : y(x_0) = y_0\}, \qquad y_0\in \mathbb{R}
\end{equation*}
Todo punto crítico $y_\ast \in D\cap C^2(I)$, además de cumplir la ecuación de Euler-Lagrange, cumple la condición extra\footnote{La estructura del problema lo exige, pues quitamos una condición y esta tiene que estar por otra parte para tener solución a una EDO de orden 2.} $F_p(x_1,y_\ast(x_1), y_\ast'(x_1)) = 0$.\\

\noindent
Esto se debe a que:
\begin{align*}
    \delta_\varphi \cc{F}(y) &=\int_{x_0}^{x_1} \left(F_y - \frac{d}{dx}(F_p)\right)\varphi~dx + \left[F_p(x,y(x),y'(x))\varphi(x)\right]_{x_0}^{x_1} \\
                             &= \int_{x_0}^{x_1} \left(F_y(x,y(x),y'(x))-\frac{d}{dx}\left[F_p(x,y(x),y'(x))\right]\right)\varphi(x)~dx \\
                             &\quad + F_p(x_1,y(x_1),y'(x_1))\varphi(x_1) - F_p(x_0,y(x_0),y'(x_0))\varphi(x_0) 
\end{align*}
En este caso, las direcciones admisibles son:
\begin{equation*}
    \cc{D} = \{\varphi \in C^1(I) : \varphi(x_0) = 0\}
\end{equation*}
Así, si nos restringimos a usar $\varphi \in C_0^1(I)\subset \cc{D}$ por la demostración del Teorema~\ref{teo:principal} recuperamos la ecuación de Euler-Lagrange para $y_\ast$:
\begin{equation*}
     F_y - \frac{d}{dx}(F_p)  = 0
\end{equation*}
Si ahora nos restringimos a todo el conjunto de direcciones admisibles $\varphi \in \cc{D}$ se sigue cumpliendo la ecuación de Euler-Lagrange pero ahora también obtenemos que $\varphi(x_0) = 0$, con lo que:
\begin{equation}\label{eq:igualdad_0}
     F_p(x_0,y(x_0),y'(x_0))\varphi(x_0)  = 0
\end{equation}
Así, si $y_\ast$ era un punto crítico, tendremos que $\delta_\varphi \cc{F}(y_\ast) = 0$, y en vista de la ecuación de Euler-Lagrange y de~\eqref{eq:igualdad_0} obtenemos que:
\begin{equation*}
     F_p(x_1,y(x_1),y'(x_1))\varphi(x_1)  = 0
\end{equation*}

\begin{ejemplo}
    Si volvemos al último ejemplo, para $I = [x_0,x_1]$, fijamos $y_0\in \mathbb{R}$, se pide ahora determinar las curvas de menor longitud que unen el punto $(x_0,y_0)$ con algún punto de abcisa $x_1$.\\

    \noindent
    Buscamos curvas descritas mediante gráficas, $x\longmapsto (x,y(x))$, obteniendo la longitud:
    \begin{equation*}
        L(y) = \int_{x_0}^{x_1} \sqrt{1+{(y'(x))}^{2}}~dx 
    \end{equation*}
    a optimizar sobre $D = \{y\in C^1([x_0,x_1]):y(x_0) = y_0\}$. Para la ecuación de Euler-Lagrange tenemos que calcular:
    \begin{equation*}
        F = F(p) = \sqrt{1+p^2} \quad\Longrightarrow\quad F_p = \frac{p}{\sqrt{1+p^2}}
    \end{equation*}
    De donde la ecuación de Euler-Lagrange quedaría:
    \begin{equation*}
        F_y - \frac{d}{dx}(F_p) = -\frac{d}{dx}\left(\frac{y'(x)}{\sqrt{1+{(y'(x))}^{2}}}\right) = 0
    \end{equation*}
    Es la misma que obteníamos en el ejemplo anterior, obteniendo que los extremales son de la forma:
    \begin{equation*}
        y(x) = c_1 + c_2x, \qquad c_1,c_2\in \mathbb{R}
    \end{equation*}
    Imponiendo las condiciones que tenemos, $y(x_0) = y_0$ y que $y$ alcanza algún punto de abcisa $x_1$ está implícito en el dominio de definición de $y$. Obligando ahora la segunda condición que sabemos que han de cumplir los puntos críticos:
    \begin{equation*}
        F_p(x_1,y(x_1),y'(x_1)) = 0
    \end{equation*}
    Tenemos que:
    \begin{equation*}
        \frac{y'(x_1)}{\sqrt{1+{(y'(x_1))}^{2}}} = 0
    \end{equation*}
    En este caso esta condición se simplifica a $y'(x_1) = 0$, con lo que $y$ ha de cumplir:
    \begin{equation*}
        \left.\begin{array}{l}
            y(x_0) = y_0 \\
            y'(x_1) = 0
        \end{array}\right\} \quad\Longleftrightarrow\quad \left\{\begin{array}{l}
            c_2 = 0 \\
            c_1 = y_0
        \end{array}\right.
    \end{equation*}
    Con lo que el único candidato a mínimo es $y(x) = y_0 \quad \forall x\in I$.
\end{ejemplo}~\\

\noindent
El resultado probado anteriormente puede generalizarse, obteiendo un resultado análogo si fijamos el extremo derecho, obteniendo la condición:
\begin{equation*}
    F_p(x_0,y_\ast(x_0), y_\ast'(x_0)) = 0
\end{equation*}
Más aún, si no se fija ningún extremo, se cumplen simultáneamente las condiciones:
\begin{equation*}
    F_p(x_0,y_\ast(x_0), y_\ast'(x_0)) = F_p(x_1,y_\ast(x_1), y_\ast'(x_1)) = 0
\end{equation*}

\begin{ejemplo}
    Como otro ejemplo, buscamos las curvas de longitud mínima que unen un punto con abscisa $x_0$ con otro con abcisa $x_1$, con $x_0<x_1$.\\

    \noindent
    Buscamos curvas descritas mediante gráficas, $x\longmapsto (x,y(x))$, obteniendo la longitud:
    \begin{equation*}
        L(y) = \int_{x_0}^{x_1} \sqrt{1+{(y'(x))}^{2}}~dx 
    \end{equation*}
    a optimizar sobre $D = \{y\in C^1([x_0,x_1])\}$. Para la ecuación de Euler-Lagrange tenemos que calcular:
    \begin{equation*}
        F = F(p) = \sqrt{1+p^2} \quad\Longrightarrow\quad F_p = \frac{p}{\sqrt{1+p^2}}
    \end{equation*}
    De donde la ecuación de Euler-Lagrange quedaría:
    \begin{equation*}
        F_y - \frac{d}{dx}(F_p) = -\frac{d}{dx}\left(\frac{y'(x)}{\sqrt{1+{(y'(x))}^{2}}}\right) = 0
    \end{equation*}
    Es la misma que obteníamos en el ejemplo anterior, obteniendo que los extremales son de la forma:
    \begin{equation*}
        y(x) = c_1 + c_2x, \qquad c_1,c_2\in \mathbb{R}
    \end{equation*}
    Ahora no tenemos condiciones en los extremos, con lo que el problema fuerza las condiciones:
    \begin{equation*}
        F_p(x_0,y(x_0),y'(x_0)) = F_p(x_1,y(x_1),y'(x_1)) = 0
    \end{equation*}
    Que en nuestro caso son:
    \begin{equation*}
        \frac{y'(x_0)}{\sqrt{1+{(y'(x_0))}^{2}}} = \frac{y'(x_1)}{\sqrt{1+{(y'(x_1))}^{2}}} = 0 \quad\Longleftrightarrow\quad y'(x_0) = 0 = y'(x_1) \quad\Longleftrightarrow\quad c_2 = 0
    \end{equation*}
    Vemos que el parámetro $c_1$ queda libre. Intuitivamente, esto se corresponde a que en cada altura tenemos un candidato a mínimo, $y(x) = c_3$.
\end{ejemplo}

\section{Solucionar Ecuaciones Diferenciales}
\noindent
Como ya ha podido percibir, la ecuación de Euler-Lagrange es una ecuación diferencial, por lo que para su uso en esta asignatura conviene adquirir los conocimientos (o repasarlos) de resolución de ecuaciones diferenciales. Un básico conocimiento de las ecuaciones diferenciales es suficiente para su desarrollo en este tema, junto con las técnicas de resolución de ecuaciones diferenciales que a continuación se describen.\\

\noindent
Resolviendo en casos prácticos la ecuación de Euler-Lagrange podemos obtener ecuaciones diferenciales del tipo:

\begin{itemize}
    \item De variables separadas.
    \item Ecuaciones lineales homogéneas con coeficientes constantes.
    \item Ecuaciones de tipo Euler homogéneas.
    \item Versiones no homogéneas de los tres casos anteriores.
\end{itemize}
Los primeros dos grandes tipos de ecuaciones diferenciales que se tratan en la enumeración anterior fueron vistos ya en la asignatura de Ecuaciones Diferenciales I, y conviene repasar el documento relativo a su temario si no se recuerda cómo solucionar este tipo de ecuaciones. Además, será adecuado que se repasen los cambios de variable de ecuaciones diferenciales.

Respecto a los dos últimos puntos de la enumereación, los tratamos de forma detallada a continuación.

\subsection{Ecuaciones de tipo Euler}
\noindent
Las ecuaciones de tipo Euler son aquellas ecuaciones diferenciales lineales en las que la variable independiente aparece multiplicada por la variable dependiente siendo la variable independiente un monomio del mismo grado que el orden de la derivada por la que se multiplica, es decir, algo del estilo:
\begin{equation*}
    x^2 y'' + 2x y' + 4y = 0
\end{equation*}
Resolver ecuaciones de tipo Euler es tan sencillo como aplicar el cambio de variable:
\begin{equation*}
    u = e^x
\end{equation*}
Y obtendremos una ecuación lineal homogénea con coeficientes constantes, que ya sabemos resolver.

\subsection{Versiones no homogéneas. Método de los coeficientes indeterminados}
\noindent
Para resolver una lineal completa se obtiene una solución particular y se le suman las soluciones de la homogénea asociada, tal y como vimos en Ecuaciones Diferenciales I.\\

\noindent
Hay que buscar clases de funciones que sean cerradas por la parte no homogénea de la ecuación diferencial:
\begin{enumerate}
    \item Si tenemos un polinomio como ecuación, buscamos con un polinomio.
    \item Si tenemos una función trigonomética, buscamos con una función trigonométrica de la misma frecuencia.
    \item Si tenemos una exponencial, buscamos con una exponencial.
\end{enumerate}

\begin{ejemplo}
    Veamos varios ejemplos:
    \begin{enumerate}
        \item  Tenemos:
            \begin{equation*}
                y''  - y = x
            \end{equation*}
            Vemos a ojo que $y_p(x) = -x$, pero en general habríamos probado con:
            \begin{equation*}
                a+bx, \quad a+bx+cx^2, \quad \ldots
            \end{equation*}
            Dejamos los coeficientes libres y los determinamos exigiendo que se cumpla la ecuación.
        \item[3.] Para la ecuación:
            \begin{equation*}
                16e^x - 8y - 2 y'' = 0
            \end{equation*}
            Probamos con $y_p(x) = ce^x$, con $c\in \mathbb{R}$:
            \begin{equation*}
                16e^x - 8ce^x - 2ce^x = 0 \quad\Longleftrightarrow\quad 16 = 10c \quad\Longleftrightarrow\quad c = \nicefrac{16}{10}
            \end{equation*}
        \item[2.] Tenemos la ecuación (donde $\alpha\in \mathbb{R}$):
            \begin{equation*}
                y'' + \alpha y = \alpha \sen(8x)
            \end{equation*}
            Probamos con $a\sen(8x) + b\cos(8x)$. Observamos que tenemos una derivada segunda y que no tenemos primera, con lo que podemos tomar ya $b = 0$:
            \begin{equation*}
                -64a\sen(8x) + \alpha a\sen(8x) = \alpha\sen(8x) \quad\Longleftrightarrow\quad (\alpha-64)a= \alpha \quad\Longleftrightarrow\quad a = \frac{\alpha}{\alpha-64}
            \end{equation*}
            Así, para $\alpha \neq 64$ tenemos dicho $a$.\\

            \noindent
            ¿Qué pasa cuando $\alpha = 64$? Resulta que $\sin(8x)$ nunca va a salir. Al igual que con raíces múltiples, si algo no funciona se pone:
            \begin{equation*}
                ax\sen(8x) + bx\cos(8x)
            \end{equation*}
            Si no funciona:
            \begin{equation*}
                ax^2\sen(8x) + bx^2\cos(8x)
            \end{equation*}
            Para $\alpha = 64$ sale $y_p(x) = -4x\cos(8x)$.
    \end{enumerate}
\end{ejemplo}


\section{Estudio de extremos locales de funcionales, condiciones suficientes}
\noindent
Dado $\cc{F}:D\to \mathbb{R}$ e $y_\ast \in D$ punto crítico, queremos estudiar su carácter de extremo local, es decir, si es o no un máximo local.

Basta estudiar el funcional en puntos de la forma $y_\varepsilon = y_\ast+\varepsilon\varphi$ para $\varphi$ dirección admisible y $\varepsilon$ pequeño ($|\varepsilon| << 1$).\\

\noindent
Desarrollando por Taylor en $\varepsilon$, para comparar $y_\varepsilon$ con $y_\ast$ se escribe:
\begin{equation*}
    \cc{F}(y_\ast + \varepsilon\varphi) = \cc{F}(y_\ast) +\varepsilon \frac{d}{d\varepsilon} \cc{F}(y_\ast + \varepsilon\varphi) \big|_{\varepsilon = 0} + \frac{\varepsilon^2}{2} \frac{d^2}{d\varepsilon^2} \cc{F}(y_\ast+\varepsilon\varphi)\big|_{\varepsilon=0} + O(\varepsilon^3)
\end{equation*}
Estamos haciendo realmente un estudio de la función de una variable
\begin{equation*}
    \varepsilon\longmapsto \cc{F}(y_\ast + \varepsilon\varphi)
\end{equation*}
Y nos interesa comparar $\cc{F}(y_\ast)$ con $\cc{F}(y_\ast + \varepsilon\varphi)$. Si tenemos un punto crítico la primera derivada será igual a 0, con lo que tendremos que estudiar el desarrollo hasta orden 2, es decir, su signo.\\

\noindent
No estudiaremos el problema en su total generalidad debido a su complejidad, por lo que sólamente nos restringiremos a un caso más sencillo, en el que tenemos un intervalo $I = [x_0,x_1]$, una función $F\in C^3(I\times \mathbb{R}\times \mathbb{R})$, la función a considerar es:
\begin{equation*}
    \cc{F}(y) = \int_{x_0}^{x_1} F(x,y(x),y'(x))~dx 
\end{equation*}
Definido sobre $D = \{y\in C^1(I) : y(x_0) = y_0, y(x_1) = y_1\}$, para ciertos $y_0,y_1\in \mathbb{R}$. Este caso ya lo tratábamos anteriormente\footnote{Salvo que ahora exigimos un mayor orden de derivabilidad, por condiciones técnicas.}. Sabemos así que las direcciones admisibles son $\varphi \in C_0^1(I)$.\\

\noindent
Como nos interesa estudiar el desarrollo de segundo orden que obtenemos al hacer ``Taylor'', dados $y\in D$, $\varphi \in C_0^1(I)$, consideramos $y_\varepsilon = y+\varepsilon\varphi$ y calculamos:
\begin{align*}
    \frac{d^2}{d\varepsilon^2}\cc{F}(y_\varepsilon) &= \frac{d}{d \varepsilon} \int_{x_0}^{x_1} \frac{d}{d\varepsilon}[F(x,y_\varepsilon(x),y_\varepsilon'(x))]~dx  \\
                                                    &= \frac{d}{d\varepsilon} \int_{x_0}^{x_1} [F_y(x,y_{\varepsilon}(x),y_\varepsilon'(x))\varphi(x) + F_p(x,y_\varepsilon(x),y_\varepsilon'(x))\varphi'(x)]~dx  \\
                                                    &= \int_{x_0}^{x_1} F_{yy} (x,y_\varepsilon(x),y_\varepsilon'(x))\varphi^2(x) + F_{yp}(x,y_\varepsilon(x),y_\varepsilon(x))\varphi(x)\varphi'(x) \\
                                                    &\quad + F_{py}(x,y_\varepsilon(x),y_\varepsilon'(x))\varphi(x)\varphi'(x) + F_{pp}(x,y_\varepsilon(x),y_\varepsilon'(x)){(\varphi'(x))}^{2}~dx 
\end{align*}
Usando ahora que $F_{yp} = F_{py}$ podemos agrupar y además agrupar $2\varphi(x)\varphi'(x)$ como:
\begin{equation*}
    2\int_{x_0}^{x_1} F_{yp}(x,y_\varepsilon(x),y_\varepsilon'(x))\varphi(x)\varphi'(x)~dx  = \int_{x_0}^{x_1} F_{yp}(x,y_\varepsilon(x),y_\varepsilon'(x)) \frac{d}{dx}{(\varphi(x))}^{2}~dx 
\end{equation*}
Integrando por partes y teniendo en cuenta que $\varphi \in C^1_0(I)$:
\begin{equation*}
     \int_{x_0}^{x_1} F_{yp}(x,y_\varepsilon(x),y_\varepsilon'(x)) \frac{d}{dx}{(\varphi(x))}^{2}~dx = -\int_{x_0}^{x_1} \frac{d}{dx}\left[F_{yp}(x,y_\varepsilon(x),y_\varepsilon'(x))\right]\varphi^2(x)~dx 
\end{equation*}
Agrupando todo esto:
\begin{align*}
    \frac{d^2}{d\varepsilon^2} \cc{F}(y_\varepsilon) &= \int_{x_0}^{x_1} F_{pp}(x,y_\varepsilon(x),y_\varepsilon'(x)){(\varphi'(x))}^{2} \\
                                                     &\quad + \left(-\frac{d}{dx}[F_{yp}(x,y_\varepsilon(x),y_\varepsilon'(x))] + F_{yy}(x,y_\varepsilon(x),y_\varepsilon'(x))\right)\varphi^2(x)~dx 
\end{align*}
Finalmente, tenemos que:
\begin{align*}
    \left[\frac{d^2}{d\varepsilon^2}\cc{F}(y)\right]_{ \big|_{\varepsilon = 0} }&= \int_{x_0}^{x_1} F_{pp}(x,y(x),y'(x)){(\varphi'(x))}^{2} \\
                                                                            &\quad + \left(-\frac{d}{dx}[F_{yp}(x,y(x),y'(x))] + F_{yy}(x,y(x),y'(x))\right)\varphi^2(x)~dx 
\end{align*}

\begin{notacion}
    Usaremos por comodidad:
    \begin{align*}
        P &= F_{pp}(x,y_\varepsilon(x),y_\varepsilon'(x)) \\
        Q &= -\frac{d}{dx}[F_{yp}(x,y_\varepsilon(x),y_\varepsilon'(x))] + F_{yy}(x,y_\varepsilon(x),y_\varepsilon'(x))
    \end{align*}
    Aunque no lo hacemos explícito, $P$ y $Q$ dependen de $x$ y de $y$.
\end{notacion}

\noindent
Y lo que nos interesa es ver el signo de
\begin{equation*}
    \left[\frac{d^2}{d\varepsilon^2}\cc{F}(y)\right]_{\big|_{\varepsilon = 0} 
}\end{equation*}
A partir de esta formulación, se puede extraer una condición necesaria para decidir si $y_\ast$ es un extremo relativo:\\

\noindent
Diremos que $y_\ast\in D$ satisface la \textbf{condición necesaria de Legendre} para mínimo local si:
\begin{equation*}
    P\geq 0 \qquad \forall x\in I
\end{equation*}
Si la condición necesaria de Legendre no se cumple entonces $y_\ast$ no nos da un mínimo local (esta afirmación requiere una prueba que saltamos por falta de tiempo). En el caso de máximo, la condición necesaria de Legendre es $P \leq 0$.\\

\noindent
En notación abreviada tenemos que:
\begin{equation}\label{eq:nosecomollamarla}
    \left[\frac{d^2}{d\varepsilon^2} \cc{F}(y_\varepsilon)\right]_{\big|_{\varepsilon= 0}} = \int P{(\varphi')}^{2} + Q\varphi^2~dx
\end{equation}
Observamos que integrando por partes (y volviendo a usar que $\varphi \in C^1_0(I)$):
\begin{equation*}
    \int_{x_0}^{x_1} (P\varphi'(x))\varphi'(x)~dx = -\int_{x_0}^{x_1} \frac{d}{dx}(P\varphi'(x))\varphi(x)~dx 
\end{equation*}
Sacando factor común en~\eqref{eq:nosecomollamarla}:
\begin{equation*}
    \left[\frac{d^2}{d\varepsilon^2}\cc{F}(y_\varepsilon)\right]_{\big|_{\varepsilon=0}} = \int_{x_0}^{x_1} \left(-\frac{d}{dx}(P\varphi') + Q\varphi\right)\varphi~dx 
\end{equation*}

\begin{definicion}[Ecuación de Jacobi]
    Dado $y\in D$ punto crítico, diremos que su \textbf{ecuación de Jacobi} asociada es:
    \begin{equation*}
        -\frac{d}{dx}(P\varphi') + Q\varphi = 0 \qquad \text{en}\quad  [x_0,x_1]
    \end{equation*}
    donde la incógnita es $\varphi \in C^2$. 
\end{definicion}

\begin{definicion}[Punto conjugado]
    Diremos que $\overline{x}\in \left]x_0,x_1\right]$ es un \textbf{punto conjugado} al $x_0$ si podemos encontrar una solución $\varphi$ de la ecuación de Jacobi tal que:
    \begin{itemize}
        \item $\varphi \not\equiv 0$.
        \item $\varphi(x_0) = 0$.
        \item $\varphi(\overline{x}) = 0$.
    \end{itemize}
\end{definicion}

\begin{teo}
    Sea $I=[x_0,x_1]$, $F\in C^3(I\times\mathbb{R}\times\mathbb{R})$, consideramos $\cc{F}:D\to \mathbb{R}$
    \begin{equation*}
        \cc{F}(y) = \int_{x_0}^{x_1} F(x,y(x),y'(x))~dx 
    \end{equation*}
    donde $D = \{y\in C^1(I) : y(x_0) = y_0, y(x_1) = y_1\}$ para $y_0,y_1\in \mathbb{R}$. Sea $y\in D$ un punto crítico que verifique la condición de Legendre estricta: $\quad P>0\quad \forall x\in I$. Entonces:
    \begin{itemize}
        \item (condición necesaria) Si $y$ es mínimo local estricto, entonces no existen puntos conjugados a $x_0$ en $\left]x_0,x_1\right[$.
        \item (condición suficiente) Si $x_0$ no tiene puntos conjugados en $\left]x_0,x_1\right]$, entonces $y$ es un mínimo local estricto (``condición de Jacobi'').
    \end{itemize}
\end{teo}

\noindent
De la demostración de este Teorema (que no haremos), se deduce el siguiente Corolario:

\begin{coro}
    Bajo estas circunstancias, si $[x_0,x_1]$ es suficientemente pequeño, se cumple la condición de Jacobi.
\end{coro}

\begin{ejemplo}
    Varios ejemplos:
    \begin{enumerate}
        \item Estudiar:
            \begin{equation*}
                \cc{F}(y) = \int_{0}^{1} {(y'(x))}^{3}~dx  
            \end{equation*}
            sobre $D = \{y\in C^1([0,1]) : y(0) = 0, y(1) = 1\}$.

            Lo primero es estudiar los puntos críticos del funcional, que los obtenemos con la ecuación de Euler-Lagrange. Previo a sus cálculos, tenemos $F = F(p) = p^3$, donde:
            \begin{equation*}
                F'(p) = 3p^2
            \end{equation*}
            Así, la ecuación de Euler-Lagrange queda:
            \begin{equation*}
                0 = -\frac{d}{dx}(F_p(x,y(x),y'(x)) = -\frac{d}{dx}\left[3{(y'(x))}^{2}\right]
            \end{equation*}
            Similar a lo que hicimos el otro día, tenemos que ${(y'(x))}^{2} = c\in \mathbb{R}$, de donde $y'(x) = d\in \mathbb{R}$, es decir, los extremales son de la forma:
            \begin{equation*}
                y(x) = c_1 + c_2x, \qquad c_1,c_2\in \mathbb{R}
            \end{equation*}
            Obligando $y\in D$ vemos que $y(x) = x$. ¿Es un máximo o mínimo local? Tenemos que calcular $P$:
            \begin{equation*}
                P = F_pp, \qquad Q = F_{yy} - \frac{d}{dx}(F_{yp})
            \end{equation*}
            Calculamos usando que $F_p(x,y,p) = 3p^2$:
            \begin{equation*}
                P = F_{pp}(x,y,p) = 6p
            \end{equation*}
            Para $y(x) = x$ tenemos que $y'(x) = 1$, y por tanto:
            \begin{equation*}
                P = F_{pp}(x, y(x), y'(x)) = 6y'(x) = 6 > 0 \qquad \forall x\in [x_0,x_1]
            \end{equation*}
            Se cumple la condición necesaria de Legendre estricta. Con vistas a aplicar el Teorema, estudiamos la ecuación de Jacobi (nuestra $F$ no depende de $y$, con lo que $Q\equiv 0$):
            \begin{equation*}
                0 = -\frac{d}{dx}(P\varphi') + Q\varphi = -\frac{d}{dx}(P\varphi') = -6\varphi'' 
            \end{equation*}
            Buscando no encontrar soluciones no nulas del tipo $\varphi(x_0) = 0$, $\varphi(\overline{x}) = 0$ para algún $\overline{x}\leq x_1$, para poder garantizar que $y$ es un mínimo local estricto.

            Solucionando la ecuación de Jacobi, la solución general es:
            \begin{equation*}
                \varphi(x) = c_1 + c_2x, \qquad c_1,c_2\in \mathbb{R}
            \end{equation*}
            Imponiendo $\varphi(x_0) = 0$ ($x_0 = 0$ y $x_1 = 1$):
            \begin{equation*}
                0 = \varphi(x_0) = c_1  \quad\Longleftrightarrow\quad c_1 = 0
            \end{equation*}
            Ahora, ¿existe $\overline{x}>0$ tal que $\varphi(\overline{x}) = 0$? Como debemos tener $c_2\neq 0$ para tener $\varphi\not\equiv 0$ no, por lo que $x_0$ no tiene puntos conjugados en $\left]x_0,x_1\right]$. 

            Así, aplicando el Teorema anterior obtenemos que $y(x) = x$ es, efectivamente, un mínimo local estricto de la función $\cc{F}$.

        \item Si tenemos ahora:
            \begin{equation*}
                \cc{F}(y) = \int_{0}^{a} {(y')}^{2}-y^2~dx , \qquad a>0
            \end{equation*}
            En $D = C^1_0([0,a])$. ¿Qué puntos críticos hay y qué podemos decir de ellos?\\

            \noindent
            Planteamos la ecuación de Euler-Lagrange, ahora para $F=F(y,p) = p^2 - y^2$:
            \begin{gather*}
                F_p = 2p, \qquad F_y = -2y \\
                F_{pp} = 2, \qquad F_{yy} = -2, \qquad F_{yp} = 0
            \end{gather*}
            Con lo que la ecuación de Euler-Lagrange quedaría:
            \begin{equation*}
                0 = F_y - \frac{d}{dx}(F_p) = -2y(x) - \frac{d}{dx}(2y'(x)) \quad\Longleftrightarrow\quad y(x) + y''(x) = 0
            \end{equation*}
            La ecuación característica es $\lm^2 + 1 = 0$, de donde $\lm = \pm i$, con lo que la solución general es:
            \begin{equation*}
                y(x) = A\sen x + B\cos x, \qquad A,B\in \mathbb{R}
            \end{equation*}
            De todas ellas, ¿cuáles están en $D$?:
            \begin{equation*}
                y(0) = 0 \quad\Longrightarrow\quad B = 0 
            \end{equation*}
            Si $y(a) = 0$, entonces $A\sen a = 0$. En este caso:
            \begin{itemize}
                \item Bien $A=0$.
                \item Bien $a = k\pi, \quad k\in \mathbb{Z}$ y no se fija el valor de $A\in \mathbb{R}$.
            \end{itemize}
            Así, en general (sin imponer nada sobre $a$), el único punto crítico es $y(x) = 0$. 

            Para todo $y\in D$ tenemos $F_{pp} = 2$, con lo que $P \equiv 2 > 0$ para todo $x\in I$, es decir, se cumple la condición de Legendre estricta para mínimo local.

            ¿Hay puntos conjugados para $x_0 = 0$? La ecuación de Jacobi es:
            \begin{equation*}
                0 = -\frac{d}{dx}(P\varphi') + Q\varphi = -2\varphi''(x) -2\varphi(x) \quad\Longleftrightarrow\quad \varphi(x) + \varphi''(x) = 0
            \end{equation*}
            Que casualmente es igual a la ecuación de Euler-Lagrange que habíamos obtenido antes, por lo que su solución general es:
            \begin{equation*}
                \varphi(x) = A\sen x + B\cos x, \qquad A,B\in \mathbb{R}
            \end{equation*}
            Si $\varphi(x_0) = 0$ entonces $B=0$. Si pedimos $\varphi(\overline{x}) = 0$, entonces
            \begin{itemize}
                \item Bien $A = 0$
                \item Bien $a=k\pi, \quad k\in \mathbb{Z}$ y no exigimos nada sobre $a$.
            \end{itemize}
            Si $a\geq \pi$, entonces podemos encontrar algún $\overline{x}\leq a$ en dicha circunstancia. Así:
            \begin{itemize}
                \item Si $a<\pi$, no hay puntos conjugados, por lo que $y(x) = 0$ es un mínimo local.
                \item Si $a>\pi$, hay punto conjugados, por lo que $y(x) = 0$ no es un mínimo local.
                \item Si $a=\pi$, no podemos afirmar nada.
            \end{itemize}
    \end{enumerate}
\end{ejemplo}
